Consistent with the channel-level ablation, region-level performance was concentrated at the frontopolar pair rather than across the frontal region as a whole (Table~\ref{tab:region_performance}; Figure~\ref{fig:region_performance}). On Raja, the two frontopolar electrodes achieved precision, recall, and $F_1$ of 0.8636, 0.8934, and 0.8657, respectively, compared with 0.7564, 0.5676, and 0.6124 across the four remaining frontal electrodes. On Cao2018, the corresponding frontopolar values were 0.7389, 0.8924, and 0.7906, whereas the two remaining frontal electrodes achieved 0.7217, 0.5405, and 0.5875. The frontal-plus-frontopolar mean $F_1$ was 0.697 compared with 0.080 across 26 non-frontal electrodes on Raja, and 0.689 compared with 0.240 across 28 non-frontal electrodes on Cao2018 (Figure~\ref{fig:exp1_region_boxplot}). The frontal versus non-frontal separation was therefore larger on Raja. This cross-dataset difference coincided with differences in montage composition, particularly for midline-or-outside electrodes, with n=3 on Raja and n=12 on Cao2018. The single-channel results retained the same frontopolar concentration (Figure~\ref{fig:exp1_single_channel}).
